\chapter*{Introduction générale}
\addcontentsline{toc}{chapter}{Introduction générale}

La transformation numérique des activités de service touche directement les agences de location de voitures. Ces structures ont besoin d'un suivi fiable des véhicules, des clients, des demandes de réservation, des locations actives et des documents produits à partir de ces opérations. En l'absence d'un système centralisé, la gestion devient vite répétitive, difficile à contrôler et source d'erreurs.

Le projet présenté dans ce rapport consiste à concevoir et réaliser une application web de gestion de location de voitures développée avec Django. L'objectif est à la fois de proposer une solution fonctionnelle et de produire un rapport académique rigoureux, cohérent avec l'implémentation réalisée et adapté à une soutenance de projet de fin d'année.

Le rapport est organisé comme suit :
\begin{itemize}
  \item le premier chapitre présente le contexte, la problématique et le périmètre du projet, ainsi que l'organisation de la démarche ;
  \item le deuxième chapitre formalise les besoins fonctionnels, les acteurs, les règles métier et les mécanismes de sécurité ;
  \item le troisième chapitre traite de la conception et de l'architecture, avec les schémas UML et le modèle de données ;
  \item le quatrième chapitre détaille la réalisation et l'ensemble des fonctionnalités implémentées ;
  \item le cinquième chapitre présente les tests, la validation et les limites identifiées ;
  \item enfin, la conclusion synthétise les apports du projet et les perspectives d'évolution.
\end{itemize}
\clearpage
